\documentclass[11pt,a4paper]{article}

\usepackage[T1]{fontenc}
\usepackage[utf8]{inputenc}
\usepackage{lmodern}
\usepackage{amsmath,amssymb}
\usepackage{geometry}
\usepackage{hyperref}
\usepackage{booktabs}
\usepackage{csvsimple}

\geometry{margin=2.7cm}

\title{Ein arithmetischer Dirac-Operator\\
zur gemeinsamen Realisierung von\\
Zeta-Nullstellen und Primzahlen}
\author{Thomas Hoffbauer}
\date{\today}

\begin{document}
\maketitle

\begin{abstract}
Wir konstruieren einen selbstadjungierten arithmetischen Dirac-Operator,
dessen Spektrum gezielt nicht-triviale Nullstellen der Riemannschen
Zetafunktion reproduziert.
Auf demselben Hilbertraum wird ein zweiter Observable eingeführt,
der Primzahlen ausgibt.
Damit erscheinen Zeta-Nullstellen und Primzahlen als komplementäre
Observablen einer einheitlichen arithmetischen Dynamik.
\end{abstract}

% ------------------------------------------------------------
\section{Einleitung}

Die Hilbert--Pólya-Idee postuliert, dass die nicht-trivialen Nullstellen
der Riemannschen Zetafunktion als Eigenwerte eines selbstadjungierten
Operators erscheinen.
Während eine kanonische unendlichdimensionale Realisierung bislang
fehlt, können endlichdimensionale Konstruktionen bereits strukturelle
Zusammenhänge zwischen Spektraltheorie und Arithmetik sichtbar machen.

Ziel dieser Arbeit ist es, eine explizite Operator-Konstruktion
vorzustellen, in der

\begin{itemize}
  \item eine vorgegebene Zeta-Nullstelle als Eigenwert erscheint,
  \item und derselbe Eigenzustand zugleich eine Primzahl als zweite
        Observable trägt.
\end{itemize}

% ------------------------------------------------------------
\section{Der arithmetische Dirac-Operator}

Sei $t_k=\Im(\rho_k)$ die imaginäre Komponente der $k$-ten
nicht-trivialen Nullstelle der Zetafunktion.
Wir betrachten einen lokalen Spektralblock
\[
\mathcal Z_k=\{t_{k-m},\dots,t_k,\dots,t_{k+m}\},
\]
und definieren auf $\mathcal H\simeq\mathbb C^{2N}$ mit $N=2m+1$ den
Dirac-Typ-Operator
\[
\mathcal D =
\begin{pmatrix}
0 & A\\
A^{*} & 0
\end{pmatrix},
\qquad \mathcal D=\mathcal D^{*},
\]
mit
\[
A=\operatorname{diag}(t_{k-m},\dots,t_{k+m})+\varepsilon R .
\]

Hierbei bezeichnet $\varepsilon>0$ eine kleine Kopplungskonstante und
$R$ eine strukturierte Kopplungsmatrix, die aus arithmetischen
Familien (Mod-$12$-Klassen) gewonnen wird.

Nach Standardresultaten der Störungstheorie gilt
\[
\sigma(\mathcal D)=\{\pm\sigma_j(A)\}, \qquad
\sigma_j(A)=t_{k+j}+O(\varepsilon),
\]
insbesondere erscheint $t_k$ als Eigenwert von $\mathcal D$.

% ------------------------------------------------------------
\section{Der Primzahl-Operator}

Auf demselben Hilbertraum definieren wir den diagonalen Operator
\[
\mathcal P=\operatorname{diag}(p_{k-m},\dots,p_k,\dots,p_{k+m}),
\]
wobei $p_n$ die $n$-te Primzahl bezeichnet.

Für einen Eigenzustand $\psi_k$ mit
\[
\mathcal D\psi_k=t_k\psi_k
\]
gilt zugleich
\[
\mathcal P\psi_k=p_k\psi_k.
\]

Der Zustand $\psi_k$ kodiert somit gleichzeitig eine
Zeta-Nullstelle und eine Primzahl.

% ------------------------------------------------------------
\section{Interpretation}

Die Konstruktion führt zu einer zweifachen Beobachtungsstruktur:

\begin{center}
\begin{tabular}{ll}
\toprule
Operator & Bedeutung\\
\midrule
$\mathcal D$ & Spektrale Dynamik (Zeta-Nullstellen)\\
$\mathcal P$ & Arithmetische Ladung (Primzahlen)\\
$\psi_k$     & Gemeinsamer Zustand\\
\bottomrule
\end{tabular}
\end{center}

In diesem Bild erscheinen
\begin{itemize}
  \item Zeta-Nullstellen als \emph{Energieniveaus},
  \item Primzahlen als \emph{Quantenzahlen}.
\end{itemize}

% ------------------------------------------------------------
\section{Numerische Realisierung}

Für $k=1,\dots,100$ wurde die beschriebene Konstruktion numerisch
implementiert.
Die folgende Tabelle zeigt die ersten 100 gemeinsamen Realisierungen
von Zeta-Nullstellen und Primzahlen.

\bigskip

% --------- Tabelle aus CSV -----------------

\begin{table}[htbp]
\centering
\small
\csvreader[
    tabular= r r r r r r r,
    table head=\toprule
    $k$ & $p_k$ & $t_k$ & $\lambda_k$ & $p_k^{(\mathrm{out})}$ & $\Delta t$ & $\Delta p$\\\midrule,
    late after line=\\,
    late after last line=\\\bottomrule
]{dirac_method_first_100.csv}{}%
{\csvcoli & \csvcolii & \csvcoliii & \csvcoliv & \csvcolv & \csvcolvi & \csvcolvii}
\caption{Erste 100 gemeinsame Realisierungen von Zeta-Nullstellen und Primzahlen
durch den arithmetischen Dirac-Operator.
$\Delta t=\lambda_k-t_k$, $\Delta p=p_k^{(\mathrm{out})}-p_k$.}
\label{tab:dirac-first-100}
\end{table}

% ------------------------------------------------------------
\section{Konsequenzen}

Die Tabelle zeigt, dass
\begin{itemize}
  \item die Eigenwerte des Dirac-Operators die Zeta-Nullstellen
        stabil approximieren,
  \item der Primzahl-Operator aus denselben Zuständen konsistent
        die zugehörigen Primzahlen ausliest.
\end{itemize}

Damit werden Zeta-Nullstellen und Primzahlen als
\emph{komplementäre Observablen} einer einheitlichen arithmetischen
Dynamik sichtbar.

% ------------------------------------------------------------
\section{Ausblick}

Mögliche nächste Schritte sind:
\begin{itemize}
  \item die Erweiterung auf größere Indizes $k$,
  \item der Übergang zu einer kanonischen unendlichdimensionalen
        Grenztheorie,
  \item sowie die Untersuchung topologischer Invarianten der
        Operatorfamilie.
\end{itemize}

% ------------------------------------------------------------
\section{Schluss}

Wir haben eine explizite operatorielle Konstruktion vorgestellt, in der
Zeta-Nullstellen und Primzahlen gemeinsam realisiert werden.
Der Ansatz liefert keinen Beweis der Riemannschen Vermutung, stellt aber
eine kohärente mechanische Interpretation der Verbindung zwischen
Spektraltheorie und Arithmetik bereit.

\end{document}
